\documentclass[11pt]{article}
\usepackage[a4paper,margin=22mm]{geometry}
\usepackage{fontspec}
\setmainfont{DejaVu Serif}
\setsansfont{DejaVu Sans}
\usepackage{microtype}
\usepackage{xcolor}
\usepackage{hyperref}
\usepackage{enumitem}
\usepackage{parskip}
\definecolor{Accent}{HTML}{971D32}
\definecolor{Muted}{HTML}{666666}
\hypersetup{colorlinks=true,urlcolor=Accent,linkcolor=Accent}
\setlist{leftmargin=1.6em,itemsep=0.3em,topsep=0.35em}
\setcounter{tocdepth}{2}
\renewcommand{\contentsname}{Contents}
\newcommand{\sectionrule}{\par\vspace{-0.5em}\noindent\textcolor{Accent}{\rule{\linewidth}{0.6pt}}\par}
\begin{document}

\begin{center}
{\sffamily\bfseries\color{Accent} TOEFL WORD OF THE DAY\par}
\vspace{8mm}
{\fontsize{42}{48}\selectfont\bfseries misconstrue\par}
\vspace{2mm}
{\Large\sffamily\color{Accent} /ˌmɪskənˈstruː/\par}
\vspace{2mm}
{\small\sffamily Local audio phrase: misconstrue\par}
{\small\sffamily Audio file: audios/misconstrue.mp3\par}
\vspace{2mm}
{\large\sffamily\color{Muted} transitive verb\par}
\vspace{5mm}
{\sffamily interpret something incorrectly\par}
\vspace{6mm}
{\large To misconstrue something is to give it the wrong meaning or interpret someone's words, actions, intentions, or evidence incorrectly.\par}
\vspace{6mm}
\begin{minipage}{0.84\textwidth}
\itshape “The findings should not be misconstrued as proof that the treatment works in every case.”
\end{minipage}\par
\vspace{10mm}
\textbf{Date:} August 24th 2026\quad\textbullet\quad \textbf{Level:} B1--mid-B2 TOEFL
\vfill
Created and compiled by \textbf{Amir Kooshky}\par
\vspace{2mm}
\href{https://t.me/KooshkyTOEFL}{Telegram Channel}\quad\textbullet\quad
\href{https://www.instagram.com/kooshkytoefl}{Instagram}\quad\textbullet\quad
\href{https://t.me/KooshkyTOEFL\_pv}{Personal Telegram}
\end{center}
\newpage
\tableofcontents
\newpage

\section{Meanings and usage}
\sectionrule
\subsection{Meaning 1: Interpret something incorrectly}
\textbf{Part of speech:} transitive verb

\textbf{IPA:} /ˌmɪskənˈstruː/

\textbf{Pronunciation phrase:} misconstrue

\textbf{Local audio file:} audios/misconstrue.mp3

\textbf{Register:} formal; common in academic, legal, political, professional, and analytical writing

\textbf{Definition:} to understand or interpret the meaning, intention, or importance of something incorrectly

\textbf{In simpler words:} understand something the wrong way

\textbf{Typical pattern:} misconstrue + something + as + something

\subsubsection{Natural examples}
\begin{enumerate}
  \item Some readers misconstrued the author's criticism as a rejection of the entire theory.
  \item The findings should not be misconstrued as proof that the treatment works in every case.
  \item She worried that her silence might be misconstrued as a lack of interest.
  \item The public initially misconstrued the announcement as a promise of immediate tax cuts.
  \item His hesitation was misconstrued as unwillingness to cooperate.
  \item Students sometimes misconstrue constructive feedback as personal criticism.
  \item The researchers warned that the correlation could easily be misconstrued as evidence of causation.
  \item Her attempt to remain neutral was misconstrued by both sides as support for their opponents.
\end{enumerate}

\subsubsection{Nuance and implication}
\textbf{Central nuance:} Misconstrue means more than simply fail to understand something. It means form a particular interpretation that is wrong.

\textbf{Usual implication:} The mistake often happens because words, behavior, evidence, or intentions are ambiguous or taken out of context.

\textbf{Compared with construe:} Construe means interpret something in a particular way and does not say whether the interpretation is correct. Misconstrue specifically means interpret it incorrectly.

\subsubsection{Grammar notes}
\textbf{How to use it:} Place the statement, action, intention, evidence, or situation being misunderstood directly after misconstrue.

\textbf{What can follow it:} The most important pattern is misconstrue something as something, where as introduces the incorrect interpretation.

\textbf{Common preposition:} Use as when stating the wrong meaning given to something, as in misconstrue silence as agreement.

\textbf{Position in a sentence:} The passive form is very common: something is, was, or could be misconstrued as something.

\subsubsection{Useful patterns}
\textbf{Pattern:} misconstrue + something + as + something

\textbf{Use:} Use this when someone gives words, actions, or evidence the wrong interpretation.

\textbf{Example:} Some employees misconstrued the new policy as a threat to their jobs.

\textbf{Pattern:} be misconstrued as + noun or -ing form

\textbf{Use:} Use the passive form when the incorrect interpretation is more important than the person making it.

\textbf{Example:} Her refusal to comment was misconstrued as an admission of guilt.

\textbf{Pattern:} should not be misconstrued as

\textbf{Use:} Use this formal pattern to warn readers against an incorrect interpretation.

\textbf{Example:} The improvement should not be misconstrued as evidence that the problem has been completely solved.

\textbf{Pattern:} could be misconstrued as

\textbf{Use:} Use this when there is a risk that something may be interpreted incorrectly.

\textbf{Example:} The informal remark could be misconstrued as an official statement.

\textbf{Pattern:} misconstrue someone's intentions

\textbf{Use:} Use this when someone incorrectly understands why another person acted in a certain way.

\textbf{Example:} Please do not misconstrue my intentions; I am trying to improve the proposal, not block it.

\subsubsection{Common wrong usage}
\paragraph{Correction 1}
\textbf{Wrong:} The students misconstrued the warning like a threat.

\textbf{Correct:} The students misconstrued the warning as a threat.

\textbf{Why:} Use misconstrue something as something, not like.

\paragraph{Correction 2}
\textbf{Wrong:} His silence was misconstrued to agreement.

\textbf{Correct:} His silence was misconstrued as agreement.

\textbf{Why:} Use as before the incorrect interpretation.

\paragraph{Correction 3}
\textbf{Wrong:} The statement can misconstrue as criticism.

\textbf{Correct:} The statement can be misconstrued as criticism.

\textbf{Why:} Use the passive form because the statement receives the interpretation.

\paragraph{Correction 4}
\textbf{Wrong:} I misconstrued what he said incorrectly.

\textbf{Correct:} I misconstrued what he said.

\textbf{Why:} Misconstrue already means interpret incorrectly, so incorrectly is usually unnecessary.

\section{Word family}
\sectionrule
\subsection{construe}
\textbf{Part of speech:} transitive verb

\textbf{IPA:} /kənˈstruː/

\textbf{Definition:} to understand or interpret the meaning of something in a particular way

\textbf{Relationship to the headword:} Misconstrue is formed from construe with the prefix mis-, meaning wrongly or incorrectly.

\textbf{Grammar pattern:} construe + something + as + something

\subsubsection{Examples}
\begin{enumerate}
  \item The court construed the law narrowly.
  \item His silence could be construed as agreement.
\end{enumerate}

\subsection{misconstrual}
\textbf{Part of speech:} countable or uncountable noun

\textbf{IPA:} /ˌmɪskənˈstruːəl/

\textbf{Definition:} an incorrect interpretation of the meaning or intention of something

\textbf{Relationship to the headword:} This noun names the incorrect interpretation produced when something is misconstrued.

\textbf{Grammar pattern:} a misconstrual of + statement, action, evidence, or intention

\textbf{Usage note:} It is formal and much less common than the verb misconstrue.

\subsubsection{Examples}
\begin{enumerate}
  \item The dispute resulted from a serious misconstrual of the original agreement.
  \item The author tried to prevent misconstrual of the study's conclusions.
\end{enumerate}

\section{Common collocations}
\sectionrule
\subsection{misconstrue as}
\textbf{Applies to:} Interpret something incorrectly

\textbf{Meaning:} incorrectly interpret something as having a particular meaning

\textbf{Pattern:} misconstrue + something + as + something

\textbf{Example:} Some readers misconstrued the author's caution as opposition to the project.

\subsection{be misconstrued as}
\textbf{Applies to:} Interpret something incorrectly

\textbf{Meaning:} be interpreted incorrectly as meaning something

\textbf{Pattern:} be misconstrued as + noun or -ing form

\textbf{Example:} The statement was misconstrued as a criticism of the entire organization.

\subsection{should not be misconstrued as}
\textbf{Applies to:} Interpret something incorrectly

\textbf{Meaning:} should not be wrongly interpreted to mean something

\textbf{Example:} The absence of evidence should not be misconstrued as proof that no effect exists.

\subsection{could be misconstrued as}
\textbf{Applies to:} Interpret something incorrectly

\textbf{Meaning:} could possibly be interpreted in the wrong way

\textbf{Example:} The decision could be misconstrued as favoritism if the criteria are not explained clearly.

\subsection{misconstrue someone's intentions}
\textbf{Applies to:} Interpret something incorrectly

\textbf{Meaning:} incorrectly understand why someone acted or spoke in a particular way

\textbf{Example:} She worried that colleagues might misconstrue her intentions.

\subsection{misconstrue a statement}
\textbf{Applies to:} Interpret something incorrectly

\textbf{Meaning:} give a statement a meaning that the speaker did not intend

\textbf{Example:} Several news reports misconstrued the scientist's statement.

\subsection{misconstrue the evidence}
\textbf{Applies to:} Interpret something incorrectly

\textbf{Meaning:} draw an incorrect interpretation from evidence

\textbf{Example:} It is easy to misconstrue the evidence if the broader context is ignored.

\subsection{easily misconstrued}
\textbf{Applies to:} Interpret something incorrectly

\textbf{Meaning:} likely to be interpreted incorrectly

\textbf{Example:} Short written messages can be easily misconstrued because they provide little information about tone.

\section{Synonyms and antonyms}
\sectionrule
\subsection{Close synonyms}
\subsubsection{misinterpret}
\textbf{Part of speech:} transitive verb

\textbf{Applies to:} Interpret something incorrectly

\textbf{Shared meaning:} Both mean understand or interpret something incorrectly.

\textbf{Difference:} Misinterpret is the general and more common word. Misconstrue is more formal and is especially common when discussing the meaning or intention of words, actions, evidence, or rules.

\textbf{Example:} The students misinterpreted the graph and reached the wrong conclusion.

\subsubsection{misunderstand}
\textbf{Part of speech:} transitive verb

\textbf{Applies to:} Interpret something incorrectly

\textbf{Shared meaning:} Both describe an incorrect understanding.

\textbf{Difference:} Misunderstand is much broader and can describe any failure to understand correctly. Misconstrue specifically suggests forming the wrong interpretation of something.

\textbf{Example:} I misunderstood the instructions and submitted the wrong document.

\subsubsection{misread}
\textbf{Part of speech:} transitive verb

\textbf{Applies to:} Interpret something incorrectly

\textbf{Shared meaning:} Both can describe giving behavior or a situation the wrong interpretation.

\textbf{Difference:} Misread can mean read written words incorrectly or, figuratively, judge a situation or person's behavior wrongly. Misconstrue is more formal and focuses on assigning the wrong meaning.

\textbf{Example:} The manager misread the employees' silence as support for the proposal.

\section{Words learners may confuse}
\sectionrule
\subsection{construe}
\textbf{Part of speech:} transitive verb

\textbf{Why learners confuse them:} Misconstrue is formed directly from construe by adding the prefix mis-, meaning wrongly.

\textbf{Key difference:} Construe means interpret something in a particular way and is neutral about whether that interpretation is correct. Misconstrue means the interpretation is wrong.

\textbf{misconstrue:} The public misconstrued the announcement as a promise of immediate action.

\textbf{construe:} The court construed the phrase as applying to all public institutions.

\subsection{misinterpret}
\textbf{Part of speech:} transitive verb

\textbf{Why learners confuse them:} The two verbs overlap strongly in meaning and can sometimes be used in the same context.

\textbf{Key difference:} Misinterpret and misconstrue are very close, but misinterpret is broader and more common. Misconstrue is especially associated with wrongly interpreting words, intentions, actions, evidence, or formal language.

\textbf{misconstrue:} His hesitation was misconstrued as opposition.

\textbf{misinterpret:} The student misinterpreted the meaning of the graph.

\subsection{misquote}
\textbf{Part of speech:} transitive verb

\textbf{Why learners confuse them:} Both can involve errors in reporting or understanding what someone said.

\textbf{Key difference:} Misconstrue means interpret the meaning of someone's words incorrectly. Misquote means report the actual words incorrectly.

\textbf{misconstrue:} The reporter misconstrued the scientist's statement as a prediction.

\textbf{misquote:} The reporter misquoted the scientist by changing several important words.

\section{Etymology}
\sectionrule
\textbf{Origin:} Misconstrue is formed from the prefix mis-, meaning wrongly or incorrectly, and construe, meaning to interpret.

\textbf{Memory link:} If construe means build an interpretation, misconstrue means build the wrong interpretation.

\section{Practice exercises}
\sectionrule
\subsection{Word family choice}
Choose the best member of the word family for each blank.

\begin{enumerate}
  \item The disagreement resulted from a serious \_\_\_\_\_ of the original statement.
  \item The court chose to \_\_\_\_\_ the rule narrowly.
  \item Readers may \_\_\_\_\_ the author's criticism as a rejection of the whole theory.
\end{enumerate}

\subsection{Correct or incorrect?}
Decide whether each sentence is correct. Correct the inaccurate ones.

\begin{enumerate}
  \item Her silence was misconstrued as agreement.
  \item The evidence should not be misconstrued like proof of causation.
  \item Some readers misconstrued the author's intentions.
  \item The statement can misconstrue as criticism.
  \item The researchers misconstrued the results incorrectly.
\end{enumerate}

\subsection{Collocation practice}
Complete each sentence with a natural collocation.

\begin{enumerate}
  \item The remark could easily be misconstrued \_\_\_\_\_ an insult.
  \item The findings should not be \_\_\_\_\_ as evidence that the problem has disappeared.
  \item Please do not misconstrue my \_\_\_\_\_; I am trying to improve the proposal, not oppose it.
  \item Without enough context, readers may easily \_\_\_\_\_ the statement.
\end{enumerate}

\subsection{Complete answer key}
\subsubsection{Word family choice}
\begin{enumerate}
  \item \textbf{misconstrual.} The disagreement resulted from a serious misconstrual of the original statement. \emph{Use the noun misconstrual for an incorrect interpretation.}
  \item \textbf{construe.} The court chose to construe the rule narrowly. \emph{Use construe when the interpretation is not necessarily incorrect.}
  \item \textbf{misconstrue.} Readers may misconstrue the author's criticism as a rejection of the whole theory. \emph{Use misconstrue when the interpretation is incorrect.}
\end{enumerate}

\subsubsection{Correct or incorrect?}
\begin{enumerate}
  \item \textbf{Correct} \emph{Be misconstrued as is a standard passive pattern.}
  \item \textbf{Incorrect}. The evidence should not be misconstrued as proof of causation. \emph{Use misconstrue something as something.}
  \item \textbf{Correct} \emph{Misconstrue someone's intentions means understand their purpose incorrectly.}
  \item \textbf{Incorrect}. The statement can be misconstrued as criticism. \emph{Use the passive form because the statement receives the interpretation.}
  \item \textbf{Incorrect}. The researchers misconstrued the results. \emph{Misconstrue already means interpret incorrectly, so incorrectly is redundant.}
\end{enumerate}

\subsubsection{Collocation practice}
\begin{enumerate}
  \item \textbf{as.} The remark could easily be misconstrued as an insult. \emph{Use misconstrue something as something.}
  \item \textbf{misconstrued.} The findings should not be misconstrued as evidence that the problem has disappeared. \emph{Should not be misconstrued as warns against an incorrect interpretation.}
  \item \textbf{intentions.} Please do not misconstrue my intentions; I am trying to improve the proposal, not oppose it. \emph{Misconstrue someone's intentions means misunderstand why they are acting in a certain way.}
  \item \textbf{misconstrue.} Without enough context, readers may easily misconstrue the statement. \emph{Misconstrue a statement means give it an incorrect meaning.}
\end{enumerate}

\vfill
\begin{center}
\small\color{Muted} Created and compiled by \textbf{Amir Kooshky}\\
\href{https://t.me/KooshkyTOEFL}{Telegram Channel}\quad\textbullet\quad
\href{https://www.instagram.com/kooshkytoefl}{Instagram}\quad\textbullet\quad
\href{https://t.me/KooshkyTOEFL\_pv}{Personal Telegram}
\end{center}
\end{document}
